\documentclass[11pt]{article}
\usepackage[margin=0.3in,top=0.25in,bottom=0.4in]{geometry}
\makeatletter
\renewcommand{\@maketitle}{\newpage\null\vskip -3.5em\begin{center}{\LARGE \@title \par}\end{center}\par\vskip 0.2em}
\makeatother
\usepackage{amsmath,amssymb}
\usepackage{hyperref}
\usepackage{booktabs}

\title{The ``HymnD'' Theory-Uncertainty Band: Construction, and the\\
Low-$p_{D^0}$ Diffractive Artifact}
\author{}
\date{}

\begin{document}
\maketitle

\section{The three uncertainty sources}

The band shown on the $D^0$ cross-section plots (\texttt{python/cross\_section.py},
\texttt{cross\_section\_sum.py}) combines three independent sources, all evaluated
at the LHAPDF fragmentation function and central BK dipole amplitude unless
stated otherwise:

\begin{description}
\item[Replica (fit) uncertainty.] The LHAPDF grid ships 101 fit replicas.
Each replica's cross section is computed exactly like the central one; the
band is the plain mean $\pm$ standard deviation across the 101 results, at
fixed $(y,p_{D^0})$ (\texttt{load\_lhapdf\_band}). This is a genuine
statistical sample, so mean/std is the right summary.
\item[BK initial-condition (BK-IC) uncertainty.] 100 samples from a
Bayesian posterior over the BK small-$x$ evolution's initial condition,
with the fragmentation function held at the central LHAPDF replica. Same
treatment: mean $\pm$ std across the 100 dipole-amplitude posterior samples
(\texttt{load\_bk\_band}).
\item[Fragmentation-scale variation.] Not a statistical sample -- a
conventional envelope from varying the fragmentation scale
$Q=C\cdot m_T$, $m_T=\sqrt{m_c^2+p_{D^0}^2}$, with $C\in\{0.5,1,2\}$ around
the central choice $C=1$. The envelope is the min/max of the three runs at
each $(y,p_{D^0})$, and its half-width $(\max-\min)/2$ is what enters the
quadrature sum below (\texttt{load\_scale\_band}).
\end{description}

\section{Combining them in quadrature}

At each $(y,p_{D^0})$, the three contributions are added in quadrature
around the plain central curve:
\[
\sigma_{\rm total} = \sqrt{\sigma_{\rm rep}^2+\sigma_{\rm BK}^2+\sigma_{\rm scale}^2},
\qquad
\text{band} = [\,\sigma_{\rm central}-\sigma_{\rm total},\ \ \sigma_{\rm central}+\sigma_{\rm total}\,]
\]
(\texttt{\_combine\_in\_quadrature}), matching \texttt{inclusive-D0-UPC}'s own
construction exactly. This treats all three sources as independent and
roughly symmetric around the central value -- a good approximation for the
two statistical samples, but only an approximation for the scale envelope,
which is bounded (a cross section can't go negative) and not necessarily
symmetric. Section~3 shows where that approximation breaks down.

\section{What goes wrong at low $p_{D^0}$ in the diffractive channel}

At small $p_{D^0}$, $m_T\approx m_c$, so the $C=0.5$ scale choice pushes
$Q=0.5\,m_T$ down toward half the charm mass -- below the charm-production
threshold built into the LHAPDF fragmentation-function grid. There, the
charm fragmentation function (and hence the cross section computed with
that scale) is \emph{exactly zero}. This is real, physical behavior of the
grid, not a bug.

Concretely, at $y=0$ (diffractive, LHAPDF central):

\begin{center}
\begin{tabular}{lrrrrr}
\toprule
$p_{D^0}$ [GeV] & central & $\sigma_{\rm rep}$ & $\sigma_{\rm BK}$ &
scale range $[\min,\max]$ & $\sigma_{\rm scale}$ (half-width) \\
\midrule
0.2 & 0.00859 & 0.00156 & 0.0000442 & $[0,\ 0.0215]$ & 0.01073 \\
0.3 & 0.01843 & 0.00291 & 0.0000993 & $[0,\ 0.0318]$ & 0.01591 \\
0.4 & 0.02990 & 0.00395 & 0.0002005 & $[0,\ 0.0416]$ & 0.02082 \\
\bottomrule
\end{tabular}
\end{center}

(all cross sections in mb/GeV). At $p_{D^0}=0.2$~GeV, the scale envelope's
lower end is the exact zero described above, and its half-width
($0.01073$) alone already exceeds the central value ($0.00859$). Combining
in quadrature,
\[
\sigma_{\rm total}=\sqrt{0.00156^2+0.0000442^2+0.01073^2}=0.01085 \ >\ \text{central},
\]
so the naive lower edge $\sigma_{\rm central}-\sigma_{\rm total}=-0.00226$
is \emph{negative} -- meaningless for a cross section, and exactly the
symmetric-uncertainty approximation from Section~2 breaking down.

This does not happen for \textbf{exclusive}, nor for the diffractive+
exclusive \textbf{sum}: at the same point, exclusive's central value is
$\approx\!1.65$~mb/GeV, roughly 200$\times$ diffractive's -- the same
absolute scale-zero effect is comparably sized in absolute terms, but tiny
relative to exclusive's much larger central value, so it never drives
exclusive's (or the sum's) lower edge negative.

\section{How it's handled on the plot}

Three approaches were tried, in order:

\begin{enumerate}
\item \emph{Clamp to a floor.} Plot
$\max(\sigma_{\rm central}-\sigma_{\rm total},\,10^{-30})$. On a log axis,
this makes the shaded band plunge in a near-vertical spike from a normal
value at neighboring $p_{D^0}$ bins down to the bottom of the visible axis
at the one affected bin -- visually alarming and easy to mistake for a
numerical error, even though the number itself isn't wrong.
\item \emph{Drop the point.} Only shade $p_{D^0}$ bins where
$\sigma_{\rm central}-\sigma_{\rm total}>0$. This removes the spike but
leaves a bare gap with \emph{no} shading at that bin -- which reads as
``zero uncertainty here,'' the opposite of the truth.
\item \emph{Interpolate/extrapolate the lower edge (current).} For bins
where $\sigma_{\rm central}-\sigma_{\rm total}\le0$, replace the lower
edge with a value inferred from the nearest \emph{valid} neighboring bins,
log-linearly: between two valid bins, log-linear interpolation; past the
edge of the valid range (as is the case here, since $p_{D^0}=0.2$~GeV is
the very first point), log-linear \emph{extrapolation} continuing the
slope of the nearest two valid bins, rather than copying the nearest
valid value flat (a flat copy was tried first and produces a visible
plateau/blunted corner instead of the sharp taper the band naturally has
elsewhere). For the table above this gives a filled lower edge of
$\approx\!5.8\times10^{-4}$ at $p_{D^0}=0.2$~GeV.
\end{enumerate}

\paragraph{Important caveat.} The interpolated/extrapolated value is
\textbf{not a measurement} -- it is a smoothing choice made purely so the
plot stays legible and continuous, by assuming the band's shape at
$p_{D^0}=0.2$ continues the trend visible at $p_{D^0}=0.3$-$0.4$. The
actual computed lower edge at that bin is negative (equivalently, the true
scale-variation envelope there touches zero), i.e.\ the real relative
uncertainty at that single point is $\gtrsim100\%$, larger than what the
smoothed curve shows. This affects only the diffractive channel, only at
the lowest one or two plotted $p_{D^0}$ bins, well outside the region of
any physical interest for the spectrum -- but should not be read as a
precise numerical result at that specific bin.

\section{Why the band is wider at low $p_{D^0}$ in general}

Section~3 covered one specific, pathological bin. But the band is visibly
wider at low $p_{D^0}$ than at high $p_{D^0}$ \emph{everywhere}, not just at
that one point -- this is a separate, broader effect worth spelling out on
its own.

\subsection{The scale-variation piece dominates the whole band}

Breaking $\sigma_{\rm total}/\sigma_{\rm central}$ down into its three
pieces at $y=0$ (LHAPDF central, exclusive and diffractive):

\begin{center}
\begin{tabular}{lrrrrrr}
\toprule
& \multicolumn{3}{c}{exclusive} & \multicolumn{3}{c}{diffractive} \\
\cmidrule(lr){2-4}\cmidrule(lr){5-7}
$p_{D^0}$ [GeV] & rel.\ total & rel.\ rep & rel.\ scale
                & rel.\ total & rel.\ rep & rel.\ scale \\
\midrule
1.0 & 0.50 & 0.037 & 0.500 & 0.50 & 0.054 & 0.500 \\
2.0 & 0.50 & 0.034 & 0.500 & 0.50 & 0.034 & 0.500 \\
4.0 & 0.29 & 0.053 & 0.286 & 0.19 & 0.031 & 0.189 \\
8.0 & 0.27 & 0.071 & 0.243 & 0.18 & 0.032 & 0.173 \\
\bottomrule
\end{tabular}
\end{center}

(rel.\ BK-IC is omitted -- it never exceeds a few percent anywhere and
plays essentially no role.) At every $p_{D^0}$, rel.\ total $\approx$
rel.\ scale: the replica and BK-IC pieces are small and nearly flat in
$p_{D^0}$, so essentially \emph{all} of the band's width, and all of its
$p_{D^0}$-dependence, comes from the fragmentation-scale-variation
envelope alone.

\subsection{Scale variation is a stand-in for missing higher orders, and
that shrinks at large scale}

Scale variation ($Q\to C\cdot Q$, $C\in\{0.5,2\}$) is not a real physical
uncertainty in the strict sense -- it is a conventional estimate of the
size of the perturbative orders this calculation does \emph{not} include
(this is leading-order in the strong coupling; a full NLO calculation would
change less under a scale variation because it explicitly includes the next
term in the series). How much the result moves under this shift is
governed by two ingredients evaluated at $Q=C\cdot m_T$, $m_T=\sqrt{m_c^2+p_{D^0}^2}$:
the running coupling $\alpha_s(Q)$, and the DGLAP evolution of the LHAPDF
fragmentation function up to scale $Q$. Both change fast with $Q$ when $Q$
is small, and flatten out (change only logarithmically) as $Q$ grows --
this is the ordinary running of $\alpha_s$ and of parton/fragmentation
evolution, not something specific to this calculation.

At low $p_{D^0}$, $m_T\approx m_c\approx1.5$~GeV: the whole $C=0.5$-to-$2$
range sits in the region of small, rapidly-varying $\alpha_s(Q)$ and steep
fragmentation-function evolution, right next to the charm-production
threshold discussed in Section~3 -- so the same factor-of-4 range in $Q$
swings the result by a large relative amount ($\sim\!50\%$ here). At high
$p_{D^0}$, $m_T$ grows well past $m_c$, the same factor-of-4 range in $Q$
now sits in a region where $\alpha_s$ runs slowly and the fragmentation
function has settled into its scaling behavior, so the relative swing
shrinks (down to $\sim\!20$-$30\%$ by $p_{D^0}=8$~GeV in the table above).

This is the standard perturbative-QCD pattern -- missing-order
(scale-variation) uncertainty is largest near a process's characteristic
low scale and shrinks at high scale -- applied here with $m_T$ (not
$p_{D^0}$ alone) as the relevant hard scale, which is why the effect is
tied to the charm mass rather than vanishing at $p_{D^0}\to0$. The Section~3
artifact is a more extreme version of exactly this same mechanism: instead
of just varying quickly, at the very lowest $p_{D^0}$ bins in the
diffractive channel the low end of the scale range ($C=0.5$) crosses
entirely below the grid's charm-production threshold and hits zero.

\section{Summary}

The HymnD band is a straightforward quadrature sum of three uncertainty
sources that is exact for the two statistical ones (replica, BK-IC) but
only approximate for the bounded scale-variation envelope. At the lowest
diffractive $p_{D^0}$ bins, that approximation fails outright (the naive
lower edge goes negative) because the scale envelope there includes a real
physical zero from the charm-mass threshold. The plot handles this by
extrapolating a smooth-looking lower edge from neighboring bins rather than
showing the literal (negative, or floor-clamped) result -- a cosmetic fix
for legibility, not a claim that the true uncertainty at those specific
bins is smaller than it actually is.

\end{document}
